\chapter{Verification and Measurement}\label{r04:ch:verification}

The code is checked at three levels: single kernels under QEMU against a scalar
reference, whole applications by a checksum of the final state, and the K1 runs
themselves, every one of which carries its checksum.

\section{Instruction sets under QEMU}\label{r04:sec:qemu}

\file{tests/test_rvv_qemu.py} in the pystencils submodule (new; a simplified version of
the 2.x test) generates small kernels with fixed-shape fields, embeds inputs and the
kernel in a C++ program, cross-compiles it with \code{riscv64-linux-gnu-g++}, runs it
under \code{qemu-riscv64 -cpu rv64,v=true,vlen=<V>}, and compares the outputs with the
same kernel compiled by the host's JIT without vectors. Result: 30 tests pass.
\begin{itemize}
  \item \isa{rvv}, \isa{rvv_ext}, \isa{-vf}, \isa{-fma}, \isa{-vf-fma}, \isa{-pf128},
    \isa{-m2} and \isa{-m4-vf-fma-pf128} at VLEN 128, 256 and 512, each from one
    \code{-march=rv64gcv} binary;
  \item \isa{-vls128} at VLEN 128 and 256; \isa{-vls256} and
    \isa{-vls256-m2-vf-fma-pf128} at 256;
  \item \isa{rvv_ext} without options prints exactly the code of \isa{rvv};
  \item \isa{rvv_ext-vls256} code does not compile with \code{zvl128b}.
\end{itemize}

\section{LBM kernels under QEMU}

\file{scripts/notes8/lbm_qemu_check.py} applies the same harness to real lbmpy
kernels on fixed $13\times6\times5$ fields; the odd $x$ extent exercises the last,
partial vector. The column is the largest difference divided by the array's largest
value, because a per-value ratio is meaningless for velocities near zero.
\begin{center}\small
\begin{tabularx}{\linewidth}{@{}>{\raggedright\arraybackslash}p{3cm}>{\raggedright\arraybackslash}p{2.2cm}>{\raggedright\arraybackslash}X>{\raggedright\arraybackslash}p{1.4cm}>{\raggedright\arraybackslash}p{2.1cm}@{}}
  \toprule
  \textbf{kernel} & \textbf{variants} & \textbf{instruction sets} & \textbf{VLEN} &
    \textbf{max diff}\\
  \midrule
  D3Q19 SRT, float64 & pull, AA even, AA odd & \isa{rvv}, \isa{rvv_ext-vf-fma},
    \isa{-vf-fma-pf128}, \isa{-vls128-vf-fma-pf128} & 128, 256 & $\leq 5.3\cdot10^{-12}$
    (per value)\\
  waLBerla-wind D3Q27 cumulant, float32 & pull & \isa{rvv}, \isa{rvv_ext-vf-fma-pf128},
    \isa{-vls256-vf-fma} & 128, 256 & $2.99$, $3.07$, $3.06\cdot10^{-5}$\\
  \bottomrule
\end{tabularx}
\end{center}

\section{The checksum}\label{r04:sec:checksum}

Both applications print, after the run, the sum of the density and the sum of
$|\mathbf{u}|^2$ over all interior cells, reduced over all processes and printed with 17
digits. A variant is compared with the reference variant (\code{pull_none}) at the same
process count and number of steps. A correct variant differs by rounding only; a wrong
kernel, remainder loop or AA even/odd step changes the sums at $10^{-3}$ or more. The
tolerance is $10^{-6}$ relative in float32 and $10^{-10}$ in float64. The sums do not
say \emph{where} an error is; that is what the QEMU checks do, cell by cell.

\section{Application checks on x86}\label{r04:sec:app-checks}

On the laptop (x86, scalar variants only) the checksums confirm that AA streaming works
through the legacy core in both applications:
\begin{center}\small
\begin{tabular}{@{}llrll@{}}
  \toprule
  \textbf{app} & \textbf{processes, domain} & \textbf{variant} & \textbf{sum($\rho$)} &
    \textbf{sum($|\mathbf{u}|^2$)}\\
  \midrule
  waLBerla-wind, float32 & 1, $64^3$ & pull & 262143.99828010798 & 757.63847218824912\\
   & & AA & 262143.99849921465 & 757.63847086051658\\
   & 8, $128^3$ & pull & 2097151.9898673892 & 7205.0347198385698\\
   & & AA & 2097151.9896876812 & 7205.034717604256\\
  \bottomrule
\end{tabular}
\end{center}
Pull and AA agree to about $10^{-9}$ relative, float32 rounding. In
\code{UniformGridCPU} (float64) they agree to $3.6\cdot10^{-16}$.

\section{The benchmark scripts}\label{r04:sec:scripts}

\file{scripts-legacy/} configures, builds, runs and tabulates both applications
(\file{scripts-legacy/README.md} has the details):
\begin{description}[style=nextline]
  \item[\file{experiments.py}] typed definitions only: machines, builds (one CMake tree
    each), variants and experiments. A variant is
    \code{Variant(streaming, vectors, vlen, lmul, vf, fma, prefetch)}, e.g.\
    \code{Variant(AA, RVV_EXT, vf=True, fma=True, prefetch=128)}.
  \item[\file{bench.py}] runs on the machine itself: \code{setup}, \code{configure},
    \code{build}, \code{run}, \code{all}, \code{list}, \code{table}. It is the only place
    that turns a \code{Variant} into a name (\code{aa_rvv_ext-vf-fma-pf128}), and before
    configuring it asks \ps{}'s \code{parse_options} whether every name means what the
    \code{Variant} says.
  \item[\file{remote.py}] runs on the laptop: pushes the whole git repository to the
    board, starts \file{bench.py} detached in \code{screen}, shows its state and fetches
    the results.
\end{description}
The timing protocol is that of the earlier reports: warm-up steps, then timed steps,
every variant once per round in turn, several rounds; one process pinned with
\code{taskset -c 2}, more with \code{mpirun --bind-to core --map-by core}; one $64^3$
block per process. \code{all} first runs a quick check (2 + 10 steps, one round) and
times only if every quick run succeeds and passes its checksum. Every run becomes one
folder \file{results-legacy/<date>-<time>-<random>/} with one JSON5 file holding the
experiment, the machine, compiler and library versions, branch, commit and uncommitted
changes of every repository, each executable's hash and flags, every single run with
its load average, and the summary.

\section{First runs on the K1}

When this report was written, the waLBerla-wind experiment had built all nine variants
on the K1 and run its quick check. Single process, 13 time steps (2 warm-up and 11
timed: the AppGenerator's \file{main.cpp} runs one step more than
\code{Parameters.timesteps}), one run each, with another user's jobs keeping the load average near 4: these numbers show
that every variant runs and is correct, but they are \emph{not} a measurement.
\begin{center}\small
\begin{tabular}{@{}lrl@{}}
  \toprule
  \textbf{variant} & \textbf{MLUP/s} & \textbf{checksum vs \texttt{pull\_none}}\\
  \midrule
  \code{pull_none} & 0.51 & (reference)\\
  \code{aa_none} & 0.57 & ok, $2.8\cdot10^{-10}$\\
  \code{pull_rvv} & 0.70 & ok, $5.4\cdot10^{-7}$\\
  \code{aa_rvv} & 0.51 & ok, $5.4\cdot10^{-7}$\\
  \code{aa_rvv_ext-vf} & 0.57 & ok, $5.5\cdot10^{-7}$\\
  \code{aa_rvv_ext-fma} & 0.52 & ok, $5.4\cdot10^{-7}$\\
  \code{aa_rvv_ext-vf-fma} & 0.62 & ok, $5.4\cdot10^{-7}$\\
  \code{aa_rvv_ext-vf-fma-pf128} & 1.07 & ok, $5.4\cdot10^{-7}$\\
  \code{aa_rvv_ext-vls256-vf-fma-pf128} & 0.67 & ok, $5.4\cdot10^{-7}$\\
  \bottomrule
\end{tabular}
\end{center}
Two observations to follow up in the timed runs:
\begin{itemize}
  \item All RVV variants differ from the scalar ones by the same amount in
    sum($|\mathbf{u}|^2$), about $5.4\cdot10^{-7}$ relative after 13 steps, while the two
    scalar variants agree to $3\cdot10^{-10}$ and the RVV variants among themselves to
    about $10^{-8}$. That is inside the tolerance, but systematic. Its source is not yet
    known.
  \item All four-process runs failed with the bus error of \cref{r04:sec:misaligned}, so
    the scripts did not start timing.
\end{itemize}
